\documentclass[12pt]{article}
\usepackage{amsmath}
\usepackage{amssymb}
\usepackage{geometry}
\geometry{margin=1in}

\title{Oblique Collisions Between Two Spheres}
\author{Logan C}
\date{July 2026}

\begin{document}
\maketitle

\section{Introduction}

Consider two spheres of masses \(m_1\) and \(m_2\), with initial velocities
\(\mathbf{u}_1\) and \(\mathbf{u}_2\). During collision, the impulse acts
along the \emph{line of centres}, which we denote by the unit vector
\(\mathbf{I}\). Only the components of velocity parallel to \(\mathbf{I}\)
change during the collision; perpendicular components remain unchanged.

\section{Velocity Decomposition}

For each sphere, decompose the velocity into components parallel and
perpendicular to the impulse direction:


\[
\mathbf{u}_i = u_{i\parallel}\mathbf{I} + \mathbf{u}_{i\perp},
\qquad i = 1,2.
\]


After the collision,


\[
\mathbf{v}_i = v_{i\parallel}\mathbf{I} + \mathbf{u}_{i\perp}.
\]



\section{Newton's Law of Restitution}

Newton's law of restitution applies only to the parallel components:


\[
e = \frac{v_{2\parallel} - v_{1\parallel}}{u_{1\parallel} - u_{2\parallel}},
\]


where \(e\) is the coefficient of restitution.

Rearranging:


\[
v_{2\parallel} - v_{1\parallel}
= -e\,(u_{2\parallel} - u_{1\parallel}).
\]



\section{Conservation of Momentum}

Linear momentum parallel to \(\mathbf{I}\) is conserved:


\[
m_1 u_{1\parallel} + m_2 u_{2\parallel}
= m_1 v_{1\parallel} + m_2 v_{2\parallel}.
\]



\section{Solving for the Final Velocities}

Solving the system consisting of restitution and momentum conservation
gives the standard 1D collision formulas:



\[
v_{1\parallel}
= \frac{m_1 u_{1\parallel} + m_2 u_{2\parallel}
      - m_2 e (u_{1\parallel} - u_{2\parallel})}
      {m_1 + m_2},
\]





\[
v_{2\parallel}
= \frac{m_1 u_{1\parallel} + m_2 u_{2\parallel}
      + m_1 e (u_{1\parallel} - u_{2\parallel})}
      {m_1 + m_2}.
\]



\section{Reconstructing the Final Velocity Vectors}

The final velocity vectors are:


\[
\mathbf{v}_1 = v_{1\parallel}\mathbf{I} + \mathbf{u}_{1\perp},
\qquad
\mathbf{v}_2 = v_{2\parallel}\mathbf{I} + \mathbf{u}_{2\perp}.
\]



\section{Computing the Impulse Direction}

Given sphere centres at \((x_1,y_1)\) and \((x_2,y_2)\), the line of centres is:


\[
\mathbf{I} = \frac{(x_2 - x_1,\; y_2 - y_1)}
                  {\sqrt{(x_2 - x_1)^2 + (y_2 - y_1)^2}}.
\]



The parallel components are:


\[
u_{i\parallel} = \mathbf{u}_i \cdot \mathbf{I},
\qquad
\mathbf{u}_{i\perp} = \mathbf{u}_i - u_{i\parallel}\mathbf{I}.
\]



\section{Collision Condition}

A collision occurs when the distance between centres satisfies:


\[
d = 2R,
\]


where \(R\) is the radius of each sphere.

\section{Conclusion}

This formulation allows efficient computation of post-collision velocities
for oblique collisions in a 2D physics engine. Only the parallel components
change, and the perpendicular components remain unaffected.

\end{document}
